\begin{longtable}{ | p{5cm} | p{7.5cm} | p{2cm} | p{1cm} | }
\caption{IO MUX 寄存器列表}\label{tab:iomux_reg_list}
\\\hline
\rowcolor{lightgray}
名称 & 描述 &  地址 & 访问 \\ \hline
 \hline
        \endfirsthead
        \hline
        \rowcolor{lightgray}
名称 & 描述 &  地址 & 访问 \\
        \hline
        \endhead
        \endfoot
\hyperref[regdesc:IOMUXPINCTRL]{IO\_MUX\_PIN\_CTRL}&外设时钟配置寄存器 & 0x3FF49000 & 读/写 \\\hline
\hyperref[regdesc:IOMUXnREG]{IO\_MUX\_GPIO36\_REG} & 管脚 GPIO36 的配置寄存器 & 0x3FF49004 &读/写\\ \hline
\hyperref[regdesc:IOMUXnREG]{IO\_MUX\_GPIO37\_REG} & 管脚 GPIO37 的配置寄存器 & 0x3FF49008 & 读/写 \\ \hline
\hyperref[regdesc:IOMUXnREG]{IO\_MUX\_GPIO38\_REG} & 管脚 GPIO38 的配置寄存器 & 0x3FF4900C & 读/写 \\ \hline
\hyperref[regdesc:IOMUXnREG]{IO\_MUX\_GPIO39\_REG} & 管脚 GPIO39 的配置寄存器 & 0x3FF49010 & 读/写 \\ \hline
\hyperref[regdesc:IOMUXnREG]{IO\_MUX\_GPIO34\_REG} & 管脚 GPIO34 的配置寄存器 & 0x3FF49014 & 读/写\\ \hline
\hyperref[regdesc:IOMUXnREG]{IO\_MUX\_GPIO35\_REG} & 管脚 GPIO35 的配置寄存器 & 0x3FF49018 & 读/写 \\ \hline
\hyperref[regdesc:IOMUXnREG]{IO\_MUX\_GPIO32\_REG} & 管脚 GPIO32 的配置寄存器 & 0x3FF4901C &读/写 \\ \hline
\hyperref[regdesc:IOMUXnREG]{IO\_MUX\_GPIO33\_REG} & 管脚 GPIO33 的配置寄存器 & 0x3FF49020 & 读/写 \\ \hline
\hyperref[regdesc:IOMUXnREG]{IO\_MUX\_GPIO25\_REG} & 管脚 GPIO25 的配置寄存器 & 0x3FF49024 & 读/写 \\ \hline 
\hyperref[regdesc:IOMUXnREG]{IO\_MUX\_GPIO26\_REG} & 管脚 GPIO26 的配置寄存器 & 0x3FF49028 & 读/写 \\ \hline
\hyperref[regdesc:IOMUXnREG]{IO\_MUX\_GPIO27\_REG} & 管脚 GPIO27 的配置寄存器 & 0x3FF4902C & 读/写 \\ \hline
\hyperref[regdesc:IOMUXnREG]{IO\_MUX\_MTMS\_REG} & 管脚 MTMS 的配置寄存器 & 0x3FF49030 & 读/写 \\ \hline
\hyperref[regdesc:IOMUXnREG]{IO\_MUX\_MTDI\_REG} & 管脚 MTDI 的配置寄存器 & 0x3FF49034 & 读/写 \\ \hline
\hyperref[regdesc:IOMUXnREG]{IO\_MUX\_MTCK\_REG} & 管脚 MTCK 的配置寄存器 & 0x3FF49038 & 读/写 \\ \hline
\hyperref[regdesc:IOMUXnREG]{IO\_MUX\_MTDO\_REG} & 管脚 MTDO 的配置寄存器 & 0x3FF4903C & 读/写 \\ \hline
\hyperref[regdesc:IOMUXnREG]{IO\_MUX\_GPIO2\_REG} & 管脚 GPIO2 的配置寄存器 & 0x3FF49040 & 读/写 \\ \hline
\hyperref[regdesc:IOMUXnREG]{IO\_MUX\_GPIO0\_REG} & 管脚 GPIO0 的配置寄存器 & 0x3FF49044 & 读/写 \\ \hline
\hyperref[regdesc:IOMUXnREG]{IO\_MUX\_GPIO4\_REG} & 管脚 GPIO4 的配置寄存器 & 0x3FF49048 & 读/写\\ \hline
\hyperref[regdesc:IOMUXnREG]{IO\_MUX\_GPIO16\_REG} & 管脚 GPIO16 的配置寄存器 & 0x3FF4904C & 读/写 \\ \hline
\hyperref[regdesc:IOMUXnREG]{IO\_MUX\_GPIO17\_REG} & 管脚 GPIO17 的配置寄存器 & 0x3FF49050 & 读/写 \\ \hline
\hyperref[regdesc:IOMUXnREG]{IO\_MUX\_SD\_DATA2\_REG} & 管脚 SD\_DATA2 的配置寄存器 & 0x3FF49054 &读/写 \\ \hline
\hyperref[regdesc:IOMUXnREG]{IO\_MUX\_SD\_DATA3\_REG} & 管脚 SD\_DATA3 的配置寄存器 & 0x3FF49058 & 读/写 \\ \hline
\hyperref[regdesc:IOMUXnREG]{IO\_MUX\_SD\_CMD\_REG} & 管脚 SD\_CMD 的配置寄存器 & 0x3FF4905C & 读/写 \\ \hline
\hyperref[regdesc:IOMUXnREG]{IO\_MUX\_SD\_CLK\_REG} & 管脚 SD\_CLK 的配置寄存器 & 0x3FF49060 & 读/写 \\ \hline
\hyperref[regdesc:IOMUXnREG]{IO\_MUX\_SD\_DATA0\_REG} & 管脚 SD\_DATA0 的配置寄存器 & 0x3FF49064 & 读/写\\ \hline
\hyperref[regdesc:IOMUXnREG]{IO\_MUX\_SD\_DATA1\_REG} & 管脚 SD\_DATA1 的配置寄存器 & 0x3FF49068 & 读/写 \\ \hline
\hyperref[regdesc:IOMUXnREG]{IO\_MUX\_GPIO5\_REG} & 管脚 GPIO5 的配置寄存器 & 0x3FF4906C & 读/写 \\ \hline
\hyperref[regdesc:IOMUXnREG]{IO\_MUX\_GPIO18\_REG} & 管脚 GPIO18 的配置寄存器 & 0x3FF49070 & 读/写 \\ \hline
\hyperref[regdesc:IOMUXnREG]{IO\_MUX\_GPIO19\_REG} & 管脚 GPIO19 的配置寄存器 & 0x3FF49074 & 读/写 \\ \hline
\hyperref[regdesc:IOMUXnREG]{IO\_MUX\_GPIO20\_REG} & 管脚 GPIO20$^1$ 的配置寄存器 & 0x3FF49078 & 读/写 \\ \hline
\hyperref[regdesc:IOMUXnREG]{IO\_MUX\_GPIO21\_REG} & 管脚 GPIO21 的配置寄存器 & 0x3FF4907C & 读/写 \\ \hline
\hyperref[regdesc:IOMUXnREG]{IO\_MUX\_GPIO22\_REG} & 管脚 GPIO22 的配置寄存器 & 0x3FF49080 & 读/写 \\ \hline
\hyperref[regdesc:IOMUXnREG]{IO\_MUX\_U0RXD\_REG} & 管脚 U0RXD 的配置寄存器 & 0x3FF49084 & 读/写 \\ \hline
\hyperref[regdesc:IOMUXnREG]{IO\_MUX\_U0TXD\_REG} & 管脚 U0TXD 的配置寄存器 & 0x3FF49088 & 读/写 \\ \hline
\hyperref[regdesc:IOMUXnREG]{IO\_MUX\_GPIO23\_REG} & 管脚 GPIO23 的配置寄存器 & 0x3FF4908C & 读/写 \\ \hline
\end{longtable}

\vspace {-1em}
\begin{small}
\begin{enumerate}
  \item GPIO20 仅在 ESP32-PICO-V3 和 ESP32-PICO-V3-02 中可用。\\
\end{enumerate}
\end{small}
